\documentclass[11pt]{article}
\usepackage[margin=1in]{geometry}
\usepackage{amsmath,amssymb}
\usepackage{mathtools}
\usepackage{booktabs}
\usepackage{microtype}
\usepackage[T1]{fontenc}
\usepackage{newtxtext,newtxmath}
\usepackage{xcolor}

\newcommand{\reac}{\mathrm{reac}}
\newcommand{\prodsup}{\mathrm{prod}}
\newcommand{\fs}{f_s}
\newcommand{\code}[1]{\texttt{\small #1}}
\DeclarePairedDelimiter{\avg}{\langle}{\rangle}

\setlength{\parskip}{0.55em}
\setlength{\parindent}{0pt}

\title{\vspace{-1.5em}\textbf{\large \texttt{Extend\_2025}: closed-form stoichiometric
mixture fraction and closure profile \(C(f)\)}}
\author{}
\date{}

\begin{document}
\maketitle
\vspace{-3.5em}
\begin{center}
\textcolor{gray!70!black}{\small Mixing-limited reaction closure \quad---\quad
\code{species\_limits.py}, \code{integrate\_species\_odes}}
\end{center}
\vspace{0.4em}

\section*{Setup}

Active species \(i\in\mathcal{A}\), reactions \(j=1,\dots,R\). Net stoichiometry matrix
\[
N\in\mathbb{R}^{|\mathcal{A}|\times R},\qquad
N_{ij}=\nu^{\prodsup}_{ij}-\nu^{\reac}_{ij}.
\]
Feed vectors \(Y^{(1)},Y^{(2)}\in\mathbb{R}^{|\mathcal{A}|}\) (stream~1 at \(f=1\),
stream~2 at \(f=0\)). The \textbf{no-reaction line} at mixture fraction \(f\in[0,1]\) is
\[
M_i(f)=f\,Y^{(1)}_i+(1-f)\,Y^{(2)}_i.
\]
Let \(y\) be the current ODE state, \(y^0\) its initial value, and \(\avg{\cdot}\) the
expectation under the step's \(\mathrm{Beta}(\alpha_t,\beta_t)\) pdf.

\texttt{Extend\_2025} shares its selectivity ray with \texttt{ray\_limit} (Section~1)
but replaces \texttt{ray\_limit}'s general breakpoint search with a \textbf{closed-form}
stoichiometric mixture-fraction formula (Section~2) --- the classical two-stream
\(f_s\) formula from reactive-mixing theory, generalized to any number of reactants per
stream.

\section*{1.\quad Selectivity ray (shared with \texttt{ray\_limit})}

The accumulated change \(\Delta:=y-y^0\) lies in \(\operatorname{range}(N)\). Since
\(\Delta=N\int_0^t r\,d\tau=N\,\xi(t)\) for the vector of per-reaction extents
\(\xi_j(t)=\int_0^t r_j\,d\tau'\ge 0\), collapsing the whole network to \emph{one}
effective reaction with \(d=\Delta\) is exact: \(\Delta\) \emph{is}
\(\sum_j\xi_j N_j\), so no extent needs to be tracked or recovered individually --- the
sum is already the ray. \texttt{Extend\_2025} reuses \texttt{ray\_limit}'s own
\code{\_selectivity\_ray} unchanged:
\[
d = \Delta = y - y^0,
\]
falling back near \(y^0\) (where \(\Delta\approx 0\)) to the mass-action projection
\(d = N\,r(y)\).

The components of \(d\) are the coefficients of the \textbf{one-line effective
reaction}: species with \(d_i<0\) are net reactants with coefficient \(\nu_i=-d_i\);
species with \(d_i>0\) are net products with coefficient \(d_i\).

\section*{2.\quad Closed-form stoichiometric mixture fraction \(\fs\)}

Rather than the general pairwise-crossing search over all consumed species
(\texttt{ray\_limit}'s \code{\_cmax\_along}), \texttt{Extend\_2025} exploits a
structural fact about two-stream feeds: if a net-reactant species is fed
\textbf{purely} by one stream, its constraint curve \(M_i(f)/\nu_i\) is a single line
through the origin at that stream's end.

\textbf{Reduction within a stream.} For reactants fed purely by stream~1
(\(Y^{(2)}_i=0\)), \(M_i(f)/\nu_i = f\cdot(Y^{(1)}_i/\nu_i)\) --- a line through
\(f=0\). Two such lines with different slopes meet only at \(f=0\), so for \(f>0\) the
one with the \textbf{smaller} slope is the minimum everywhere on \((0,1]\): the whole
group collapses to a single line with slope
\[
\kappa_1 = \min_{i\;\text{pure stream 1}}\frac{Y^{(1)}_i}{\nu_i}.
\]
Symmetrically, reactants fed purely by stream~2 (\(Y^{(1)}_j=0\)) collapse to
\[
\kappa_2 = \min_{j\;\text{pure stream 2}}\frac{Y^{(2)}_j}{\nu_j}.
\]

\textbf{The crossing.} The achievable extent is
\(c_{\max}(f)=\min\bigl(f\kappa_1,\,(1-f)\kappa_2\bigr)\) --- one line increasing from
\(0\), one decreasing to \(0\) --- so they cross at exactly one interior point:
\[
\boxed{\fs = \frac{\kappa_2}{\kappa_1+\kappa_2}}\,,\qquad
c^{\star}=\fs\,\kappa_1=(1-\fs)\,\kappa_2.
\]
This is the standard two-stream stoichiometric mixture fraction, generalized beyond a
single reactant per stream by taking the binding (minimum-ratio) reactant on each side.
The complete-reaction limit is the same \textbf{two-segment tent} as
\texttt{ray\_limit}'s (only its derivation differs):
\[
B_i(0)=Y^{(2)}_i,\qquad B_i(\fs)=M_i(\fs)+c^{\star}d_i,\qquad B_i(1)=Y^{(1)}_i.
\]

\textbf{Scope.} The reduction requires every net-reactant species to be fed by exactly
one stream. Species absent from both feeds (\(y^0_i=0\)) can never violate this ---
since concentrations are non-negative, \(d_i=y_i-y^0_i=y_i\ge 0\), so an unfed
intermediate is never a net reactant. If an \emph{actual feed} species is present in
both streams and turns out to be a net reactant, the closed form does not apply;
\texttt{Extend\_2025} does not fall back to \texttt{ray\_limit}'s general search in that
case --- it prints a diagnostic message and refuses to compute a limit for that step
(rates set to \(0\)).

\section*{3.\quad Closure profile \(C_i(f)\) and scalar \(\lambda\)}

Identical machinery to \texttt{ray\_limit} (same boxed formula, same \(\lambda\), same
segment integrals --- only \((\fs,B)\)'s construction differs):
\[
\boxed{\;C_i(f)=B_i(f)+\bigl[M_i(f)-B_i(f)\bigr]\,\lambda_i\;}
\]
\[
\lambda_i=\operatorname{clip}_{[0,1]}\frac{y_i-\avg{B_i}}{\avg{M_i}-\avg{B_i}}.
\]
Segment integrals use the regularised incomplete Beta function \(I_x(\alpha,\beta)\)
(\code{betainc}):
\[
\int_{f_a}^{f_b}(sf+b)\,p_\beta\,df
= s\,\frac{\alpha_t}{\alpha_t+\beta_t}
  \bigl[I_{f_b}(\alpha_t{+}1,\beta_t)-I_{f_a}(\alpha_t{+}1,\beta_t)\bigr]
+ b\,\bigl[I_{f_b}(\alpha_t,\beta_t)-I_{f_a}(\alpha_t,\beta_t)\bigr].
\]
Because \(B_i-M_i=c_{\max}(f)\,d_i\) throughout (same as \texttt{ray\_limit}), the same
\textbf{global \(\lambda\)} identity holds: \(\lambda_i=1-1/\avg{c_{\max}}\) is the same
scalar for every species, and \(\avg{C_i}=y_i\) exactly whenever the clamp is not
binding.

\section*{4.\quad Worked illustration: \texttt{input\_fuller\_BC\_azo\_coupling\_kinetics}}

\textbf{Scheme:} R1: \(A{+}B\!\to\!R\), R2: \(A{+}R\!\to\!S\), R3: \(A{+}B\!\to\!T\),
R4: \(A{+}T\!\to\!S\), R5: \(A{+}C\!\to\!Q\).
\textbf{Feeds:} \(Y^{(1)}=15\,A\) (stream~1), \(Y^{(2)}=1.2\,B+1.2\,C\) (stream~2);
\(f_{\mathrm{mean}}=0.0625\).

Same output-step time-midpoint as \texttt{ray\_limit}'s own worked example
(\(t=t_{\mathrm{end}}/2=0.07629\) s \(\approx 3.08\,\tau_s\), \(\tau_s=0.02476\) s,
\((\alpha_t,\beta_t)=(0.372,\,5.575)\) --- this instant is \(\approx 70\%\) \(A\)
conversion, not \(50\%\)), the ray gives the effective reaction
\[
0.6556\,A + 0.3314\,B + 0.2868\,C \;\longrightarrow\;
0.2832\,R + 0.0108\,T + 0.0374\,S + 0.2868\,Q.
\]

\textbf{Closed-form \(\fs\).} \(A\) is pure stream~1 (\(Y^{(1)}_A=15\)); \(B,C\) are
pure stream~2 (\(Y^{(2)}_B=Y^{(2)}_C=1.2\)):
\[
\kappa_1=\frac{15}{0.6556}=22.88,\qquad
\kappa_2=\min\!\left(\frac{1.2}{0.3314},\,\frac{1.2}{0.2868}\right)
=\min(3.622,\,4.184)=3.622\ (B\ \text{binds}).
\]
\[
\fs=\frac{3.622}{22.88+3.622}=0.1366,\qquad c^{\star}=\fs\,\kappa_1=3.127.
\]
Peak values \(v^{\mathrm{fs}}_i=M_i(\fs)+c^{\star}d_i\) (both \(A\) and \(B\) vanish
exactly at \(\fs\), by construction) and the resulting averages and closure scalars:

\begin{center}\small
\begin{tabular}{c|cccccc}
\toprule
Species & \(\avg{M_i}=y^0_i\) & \(v^{\mathrm{fs}}_i\) (mol/L) &
\(\avg{B_i}\) (mol/L) & \(y_i\) (mol/L) & \(\lambda_i\) & \(\avg{C_i}\) \\\midrule
\(A\) & 0.9375 & 0.000 & 0.2814 & 0.2819 & 0.000839 & 0.2819 \\
\(B\) & 1.1250 & 0.000 & 0.7934 & 0.7936 & 0.000839 & 0.7936 \\
\(C\) & 1.1250 & 0.139 & 0.8379 & 0.8382 & 0.000839 & 0.8382 \\
\(R\) & 0.0000 & 0.886 & 0.2835 & 0.2832 & 0.000839 & 0.2832 \\
\(T\) & 0.0000 & 0.034 & 0.01076 & 0.01075 & 0.000839 & 0.01075 \\
\(S\) & 0.0000 & 0.117 & 0.03740 & 0.03737 & 0.000839 & 0.03737 \\
\(Q\) & 0.0000 & 0.897 & 0.2871 & 0.2868 & 0.000839 & 0.2868 \\
\bottomrule
\end{tabular}
\end{center}

Every \(\lambda_i\) equals the same global scalar and \(\avg{C_i}=y_i\) to the quoted
digits --- identical structure to \texttt{ray\_limit}. At this same instant
\texttt{ray\_limit}'s own worked example gives \(d_A=-0.6556\), \(d_B=-0.3314\),
\(d_C=-0.2868\), \(\fs=0.1366\), \(c^{\star}=3.127\) --- matching to the quoted digits,
as expected: \(A\), \(B\), \(C\) are each fed by exactly one stream throughout this
network (and \(R,T,S,Q\), being unfed, can never violate that condition), so the closed
form applies at every step and \texttt{Extend\_2025} reproduces \texttt{ray\_limit}'s
trajectory numerically, at matching CPU cost, without ever performing the general
breakpoint search.

\section*{Note: relation to \texttt{ray\_limit}}

\texttt{Extend\_2025} and \texttt{ray\_limit} share the selectivity ray \(d\) and the
entire closure-profile machinery (Section~3); they differ only in how \((\fs, B)\) is
constructed from \(d\): \texttt{ray\_limit}'s search handles \emph{any} feed
distribution (including a reactant fed by both streams) via an
\(O(|\mathcal{C}|^2)\) pairwise-crossing search; \texttt{Extend\_2025}'s closed form
handles only the pure single-stream-per-reactant case, in closed form, and refuses
outright --- rather than silently falling back --- when that condition fails. Whenever
the condition holds (as it does for every network exercised so far), the two
constructions are mathematically identical, not merely similar: both reduce to the
unique crossing of one increasing and one decreasing affine function of \(f\).

\end{document}
